\documentclass[11pt]{article}
\usepackage[margin=1in]{geometry}
\usepackage[utf8]{inputenc}
\usepackage[T1]{fontenc}
\usepackage{enumitem}
\usepackage{hyperref}
\usepackage{xcolor}
\usepackage{titlesec}

\titlespacing*{\section}{0pt}{8pt}{4pt}
\titlespacing*{\subsection}{0pt}{6pt}{3pt}
\setlength{\parskip}{4pt}
\setlength{\parindent}{0pt}
\pagestyle{empty}

\begin{document}

\begin{center}
{\Large\bfseries Hypothesis Generation from Multi-Modal Data}\\[4pt]
{\normalsize Alvaro Velasquez, University of Colorado Boulder}\\[2pt]
{\small April 2026}
\end{center}

\vspace{-4pt}
\hrule
\vspace{6pt}

\section*{Section 1: Background and Introduction}

Consider an AI capable of making a discovery worthy of a Nobel Prize. Such transformational scientific discoveries do not simply extend the prevailing scientific assumptions and knowledge of their time; they often challenge or even contradict them. For example, the discovery of quantum mechanics contradicted Newtonian physics at the microscopic scale. The AI analogy to this is that the AI for Science community would like to train an AI on the existing research literature such that it generalizes so well that it can discover a new scientific result that invalidates its own training data. While this seems like a paradox, logicians have developed defeasible logics to handle such contradictions. We seek to integrate such symbolic logical systems with modern scientific foundation models to develop neurosymbolic defeasible inference and training methods. that combine (i) AI4Math and AI4Science models as probabilistic and hallucinatory sources of knowledge with (ii) the aforementioned logical formalisms, e.g., RuleLog, for defeasible knowledge representation and reasoning. Our goal is to develop a defeasible AI that makes a transformational scientific discovery that invalidates its own training data and is rigorously proved as a theorem. To that end, we will develop defeasible inference and training  methods that combine (i) scientific foundation models as probabilistic and hallucinatory sources of knowledge with (ii) formal defeasible logics like RuleLog for defeasible knowledge representation and reasoning.

A critical dimension of this challenge is that modern scientific data is inherently multimodal: protein structures are resolved by cryo-EM imaging, molecular dynamics are observed through spectroscopic time series, and drug--target interactions are characterized by 3D binding-site geometries alongside textual literature. The grounding problem---tracing a model's conclusions back to their supporting evidence---is fundamentally harder in multimodal settings. In text-only reasoning, entity identification is given directly (the name appears in the prompt). When evidence arrives as an image or a 3D structure, the model must first \emph{perceive} the entity before it can reason about its properties. Convergent evidence from multiple research communities establishes that this perception step introduces a systematic bias against precisely the entities that matter for defeasible reasoning:

\begin{itemize}[nosep, leftmargin=1.5em]
\item \textbf{Recognition is biased against atypical instances.} Even on balanced datasets like ImageNet, per-class accuracy ranges from 100\% to 16\% (Cui et al., CVPR 2024), and the disparity arises from learned representations, not classifier bias. Atypical class members---the entities governed by defeater rules in defeasible logic---fall on the hard end of this distribution.

\item \textbf{Multimodal metrics encode prototypicality bias.} The ProtoBias benchmark (2026) demonstrates that CLIPScore, PickScore, and VQA-based evaluation metrics systematically prefer prototypical-but-semantically-wrong images over atypical-but-correct ones. The very infrastructure used to evaluate multimodal systems is biased against exceptions.

\item \textbf{Distribution-shift robustness is illusory.} CLIP's apparent robustness to distribution shift is largely explained by in-domain examples in web-scale pretraining data; on strictly out-of-domain inputs, performance degrades substantially (ICLR 2025). Atypical instances are by definition out-of-distribution within their class.

\item \textbf{Visual abductive reasoning shows large human--VLM gaps.} On the Black Swan benchmark (CVPR 2025), which evaluates abductive and defeasible video reasoning about unexpected events, VLMs trail human performance by up to 32 percentage points.

\item \textbf{The bottleneck is integration, not perception alone.} Recent work (ICLR 2026) identifies a task-composition bottleneck in multimodal models: recognition and reasoning can each succeed in isolation but fail when composed, and early fusion introduces systematic bias---directly predicting that entity-identification-then-defeasible-reasoning chains will compound errors.
\end{itemize}

\noindent For scientific discovery, these findings create a double bottleneck: the novel exceptions a defeasible AI must find (e.g., an intrinsically disordered protein that violates the structure--function default) are also the hardest to identify from perceptual data such as cryo-EM density maps and molecular surface visualizations. Our preliminary evaluation of vision--language models on defeasible reasoning tasks in the DeFAb benchmark confirms this compounding effect empirically.

\section*{Section 2: Project Objective}

Our goal is to develop a defeasible AI that makes a transformational scientific discovery capable of invalidating its own training data. Today, foundation models are routinely evaluated on forward, monotonic inference, yet scientific progress depends on abduction and revision captured by non-monotonic, or defeasible, reasoning: proposing hypotheses that explain observations and retracting default conclusions when exceptions arise. To show these limitations, we have run defeasible reasoning tasks on state-of-the-art models GPT-5.2 and Claude Sonnet 4.6, yielding accuracy rates of 1.4\% and 2.1\%. These experiments demonstrate that current technology is not capable of executing basic defeasible reasoning, much less the complex defeasible reasoning required for making transformational scientific discoveries.

Crucially, this capability gap widens when scientific evidence is multimodal. When defeasible claims derived from textual literature must be reconciled with structural or imaging data---the setting in which most wet-lab discoveries actually occur---the model must ground its formal reasoning in perceptual input, a capability no current system reliably provides. We therefore propose to develop defeasible inference, fine-tuning, and interpretability methods that operate natively across text, molecular structure, and image modalities, maintaining traceable and revisable knowledge representations throughout.

\section*{Section 3: Proposed Research and Methods}

We propose a novel three-pronged approach detailed below:

\begin{itemize}[leftmargin=1.5em]
\item \textbf{Defeasible inference (Months 1--3):} we propose to learn a defeasible knowledge structure in RuleLog that marks which rules are strict versus revisable, traces every conclusion to its support set, and enables targeted revision of knowledge. This defeasible structure will be used to guide the discovery of prompts for defeasible inference and to retrieve supporting and invalidating information using novel defeasible retrieval-augmented generation (RAG) methods. In particular, these RAG methods, which normally retrieve supporting information for a given prompt by searching for similar vector representations of other text, will be extended to instead retrieve evidence that invalidates a claim.

We will further extend this defeasible RAG framework to multimodal retrieval. The defeasible knowledge structure will map each entity to representations across modalities---text, 3D structure, and image---so that defeater retrieval can query across modalities. For example, when a sequence-based rule predicts that a protein has a stable fold, the retrieval system can surface a cryo-EM density map or an AlphaFold confidence plot that contradicts the prediction, providing cross-modal evidence for a structural defeater. This multimodal extension directly addresses the grounding problem in multimodal embeddings: defeasible RAG must not only find semantically similar text, but identify perceptual evidence in a different modality that invalidates the textual claim.

\item \textbf{Defeasible fine-tuning (Months 4--6):} we propose to integrate reinforcement learning (RL), planning, and preference optimization into a defeasible fine-tuning framework for updating the weights of the scientific foundation model to achieve defeasible generalization. While preference optimization and RL are common approaches to fine-tuning AI models, they often rely on hand-crafted preference datasets that do not capture defeasibility and are costly to generate. We instead propose defeasibility games as a game the AI plays against itself to try to invalidate claims and find so-called defeaters that reconcile the resulting contradictions. These games will be logically formalized using the Relational Dynamic Influence Diagram Language (RDDL) and can be played using planning approaches combined with RL for more efficient search. Novel preference optimization techniques will then be deployed to fine-tune the model based on the results of the games.

To address the multimodal grounding deficit, the defeasibility games will incorporate multimodal inputs: the model plays against itself using mixed text, structure, and image evidence, forcing it to ground defeasible reasoning in perceptual data. The reward signal exploits a key property of our framework---because the polynomial-time proof checker verifies correctness at the level of the formal defeasible theory, it scores reasoning results identically regardless of input modality. This means the same verifier-backed reward function that drives text-only defeasibility games extends without modification to multimodal settings, providing an exact training signal for cross-modal defeasible generalization without additional annotation.

\item \textbf{Mechanistic interpretability for explainability and model editing (Months 7--9):} RuleLog justifications are derived from our learned defeasible knowledge structure by searching for derivations that explain the inference result, e.g., a variable may have been set to false because its associated rule was defeated by another rule or another rule was prioritized. This justification is used in tandem with the parameter updates from Thrust 2 to localize which parameters in the model capture the concepts of the justification. These parameters can be automatically edited to match the derived justification or manually edited by end-users.

In the multimodal setting, this interpretability framework extends to cross-modal concept localization: identifying which parameters encode visual representations of entities (e.g., the structural signature of a disordered region in a cryo-EM encoder) versus textual representations (e.g., the sequence-level rule ``protein $\Rightarrow$ has\_stable\_3d''), and tracing how these interact during defeasible inference. When a defeater is discovered through structural imaging data, the justification trace reveals whether the model grounded its revision in the visual evidence or in a textual proxy, enabling targeted editing of the parameters responsible for cross-modal grounding failures.
\end{itemize}

We believe this approach will be successful as RuleLog has polynomial-time inference complexity, is verifiable for our defeasible reasoning, adds a capability scientific AI models currently lack in finding defeaters for defeasible reasoning, and the defeasible knowledge structure will enable more efficient inference and fine-tuning from little data as the proposed defeasible games generate high-quality data.

\section*{Section 4: Milestones}

\begin{itemize}[leftmargin=1.5em, nosep]
\item By the end of the third month, we will discover a novel peptide (small protein $<$20 amino acids). Novelty is measured as less than 20\% sequence similarity to the training data.
\item By the end of the sixth month, we will discover a novel polypeptide ($>$50 amino acids).
\item By the end of the ninth month, we will scale to much larger proteins ($>$200 amino acids) to discover novel in silico candidates for one of the following four holy grails: polymers that prevent damaging ice crystals from forming on frozen organs, enzymes that depolymerize hard plastics like polyethylene, bacteria for self-healing space materials, or alternatives to lithium for novel batteries that mitigate our battery supply chain challenges.
\item By the end of the ninth month, we will additionally demonstrate that multimodal defeasible reasoning---combining sequence data with structural and imaging data---discovers protein candidates that text-only methods miss. Specifically, we will show that for cases where the relevant exception (defeater) is structure-dependent (e.g., disordered regions, allosteric binding sites, non-canonical folds), the multimodal pipeline identifies candidates with higher binding affinity or greater structural novelty than the text-only pipeline, validating that cross-modal grounding resolves a genuine bottleneck in defeasible scientific discovery.
\end{itemize}

\section*{Section 5: Data Sources and Models}

For drug discovery problems, we will leverage DeepPurpose for efficient binding affinity predictions between the candidate protein and the target binding site and we will use BioEmu for AI-accelerated prediction of molecular dynamics simulations. We will compare our discovered drugs against the NVIDIA GenMol drug discovery model.

For multimodal grounding, we will integrate three additional data sources already connected via our image harvesting pipeline: the Protein Data Bank (PDB), which provides experimentally determined 3D structures for over 220,000 macromolecules; the Electron Microscopy Data Bank (EMDB), which provides cryo-EM density maps that serve as the primary visual evidence for structural exceptions; and AlphaFold DB, which provides predicted structures with per-residue confidence scores that can signal structurally atypical regions. Entity-level image associations from Wikidata (P18 image property, queryable via SPARQL) and VisualSem (938,000 images across 90,000 entities bridging WordNet and BabelNet) provide cross-modal entity representations for the broader defeasible knowledge structure.

\section*{Section 6: Metrics}

We propose a new metric for generalization of scientific AI models. Conventional out-of-distribution (OOD) generalization is measured by ensuring a large KL-divergence value between the training and test distributions and then measuring accuracy on the test distribution after training the model on the training distribution. However, defeasible generalization will require the measurement of whether a generated result invalidates the scientific assumptions in its training data, which is not captured by conventional OOD metrics. Our go/no-go decision will be whether our methods discover a drug with greater binding affinity than what the NVIDIA GenMol model can produce while also exhibiting at most 20\% sequence similarity to the training set.

We additionally propose a \emph{cross-modal defeasible accuracy} metric that measures whether the model can identify defeaters when evidence arrives in a different modality from the original claim. Concretely: given a defeasible rule derived from textual literature (e.g., ``proteins with motif $X$ bind target $Y$''), the model is presented with structural or imaging evidence that constitutes an exception (e.g., a cryo-EM map showing a conformational state that blocks binding). Cross-modal defeasible accuracy is the fraction of such instances in which the model correctly identifies the defeater and revises its conclusion while preserving unrelated expectations. This metric directly operationalizes the grounding problem in multimodal embeddings and provides a measurable target for the multimodal extensions proposed in Thrusts 1--3.

\end{document}
